\begin{abstract}
Frequency estimation latency and robustness are critical for protection in
low-inertia, Inverter-Based Resource (IBR)-dominated grids. A stress-oriented
benchmark evaluates twelve estimators from five algorithmic families under five
scenarios using dual-rate simulation (1\,MHz physics / 10\,kHz DSP). Each
estimator is independently tuned per scenario via grid search, establishing the
performance ceiling per estimator--scenario pair, and assessed on RMSE, peak
error, trip-risk duration ($T_{\mathrm{trip}}$), and computational cost
($N\!=\!30$ Monte Carlo). Scenarios~(A)--(C) follow the IEC/IEEE~60255-118-1
test philosophy; Scenarios~(D)--(E) introduce composite IBR-specific disturbances
beyond the standard scope. Several estimators that pass standard compliance tests
degrade substantially under composite IBR disturbances; no single family dominates
across all scenario types. Window-based and loop-based methods exhibit
complementary strengths; data-driven approaches incur higher computational cost
without consistent accuracy gains over scenario-optimized classical methods. The
results indicate that the IEC/IEEE~60255-118-1 test scope does not fully
characterize estimator resilience in IBR-dominated conditions, and that composite
phase-step and harmonic distortion scenarios expose failure modes absent from
isolated single-disturbance tests.
\end{abstract}